\documentclass[10pt,a4paper]{amsart}
\usepackage[margin=19mm]{geometry}
\usepackage{amsmath,amssymb,mathtools}
\usepackage[scr=rsfs,cal=euler]{mathalfa}
\usepackage{xfrac}
\usepackage[unicode,hidelinks,bookmarks=false]{hyperref}
\newcommand{\ie}{i.e.\ }
\newcommand{\cf}{cf.\ }
\newcommand{\resp}{respectively\ }
\newcommand{\Subsection}{\subsection}
\providecommand{\nobreakdash}{\nobreak\hbox{-}}
\setlength{\emergencystretch}{2em}
\multlinegap0pt
\usepackage{bm}
\usepackage{bbm}
\usepackage{dsfont}
\usepackage{xspace}
\usepackage{cancel}
\usepackage{mathtools}
\usepackage{amsthm}
\makeatletter
\@ifundefined{Theorem}{\newtheorem{Theorem}{Theorem}}{}
\@ifundefined{Lemma}{\newtheorem{Lemma}[Theorem]{Lemma}}{}
\makeatother
\newcommand{\boldeta}{{\boldsymbol{\eta}}}
\newcommand{\oddd}{\mathrm{ODiff}}
\newcommand{\odd}{\mathrm{Odd}}
\title[The peak algebra in noncommuting variables]{The peak algebra in noncommuting variables}
\author{Farid Aliniaeifard, Shu Xiao Li}
\date{Source: arXiv:2506.12868v3 (CC BY 4.0)}
\begin{document}
\maketitle

\noindent\emph{Source and licence.} The peak algebra in noncommuting variables. Version arXiv:2506.12868v3 (CC BY 4.0), distributed under the Creative Commons Attribution 4.0 International licence, as recorded in the arXiv licence metadata for this exact version and shown on the abstract page https://arxiv.org/abs/2506.12868v3. Full rights notice: licenses/arxiv-2506.12868-CC-BY-4.0.txt.\par\smallskip

\noindent\textit{Editorial scope.} Counted unit: the Lemma whose source label is \texttt{lem:M} in the article (source0.tex lines 2710--2718, source label \texttt{lem:M}), delivered together with the article's own complete original proof of it (source0.tex lines 2721--2784, one \texttt{proof} environment of 64 lines). No mathematics is written, repaired or re-derived: statement and proof are the article's own text line for line. Required context retained uncounted: none. Every definition, display and label of those lines is retained unchanged. Omitted from this package and not redistributed: 1--2709, 2719--2720, 2785--4501. Nothing inside a retained excerpt is omitted or altered: every retained body is delivered exactly as the article's lines are written, so no omission receipt of any kind accompanies this package. Citations: the retained excerpts carry no citation command at all, so no bibliography entry of the article is transcribed and no reference is redistributed. Reference adaptation: none was needed and none was made; every reference inside the delivered unit resolves inside it, so no unresolved reference or citation is produced. Printed slips kept literal: no printed word, symbol, hypothesis or number of the counted statement or of its proof was altered. Layout and class: the article's own class is \texttt{amsart} with packages including geometry, comment, hyperref, amsthm, amsmath, graphicx, amsfonts, amssymb, tikz, mathrsfs, caption, stmaryrd, dsfont, marginnote; the wrapper is the lane's pinned \texttt{amsart} editorial wrapper at 10pt on A4, which restates the theorem-like environments the retained text needs and adds the article's own macro definitions that the retained text needs (\texttt{\textbackslash boldeta}, \texttt{\textbackslash odd}, \texttt{\textbackslash oddd}), with extra wrapper packages (bm, bbm, dsfont, xspace, cancel, mathtools). The delivered numbering is the unit's own. Assets: no retained span contains a figure environment, a graphics-inclusion command, a PDF-page inclusion, a table or a TikZ picture; no image file is copied into this package, the package declares no asset dependency and no image file is redistributed with it. Licence: version arXiv:2506.12868v3 of this article is distributed under the Creative Commons Attribution 4.0 International licence, as recorded in the arXiv licence metadata for this exact version and shown on the abstract page https://arxiv.org/abs/2506.12868v3. A full rights notice is delivered at licenses/arxiv-2506.12868-CC-BY-4.0.txt. Characters: every retained excerpt is pure ASCII, so the delivered engine, XeLaTeX with its default Latin Modern text font, reports no missing character.
\begin{Lemma}\label{lem:M}
	Let $A\subseteq[n-1]$ and $\sigma\in\mathfrak{S}_n$. Then
	$$\sum_{A\subseteq C\subseteq[n-1]}(-1)^{|C|}\mathbf{K}_{(C\setminus((C+1)\cup\{1\}),\sigma)}=\begin{cases}
		(-1)^{n-1-|B|}\boldeta_{(B,\sigma)} & \text{ if }n-\max (A)\text{ is odd,}\\
		0 & \text{ otherwise,}
	\end{cases}$$
	where $B$ is the unique peak set such that $\odd(B)=\oddd(A)$.
	
	\end{Lemma}

\begin{proof}

%\begin{enumerate} 
%\item Let $D$ be the set of all elements of $C\setminus \{1\}$ except if some consecutive integers $i, i+1, i+2, \dots$ are in $C$ only the smallest one $i$ is in $D$.  For example, if $C=\{2,3,4,6,7\}$, then $D=\{2,6\}$.\\
%\item  Let $D_{> \max{A}}$ be the set of all elements of $D$ that are bigger than or equal to $\max(A)$. Continuing our example, if $A=\{3,4\}$, then $D_{\geq \max{A}}=\{6\}$.\\
%\item Order the elements of $D_{\geq \max{A}}\cup \{\max(A),n+1\}$ and find its smallest element $a$ such that its difference with the next element $b$ is odd. Continuing our example, $D_{\geq \max{A}}\cup \{\max(A),n+1\}=\{4,6,9\},$ $a=6$, $b=9$.
%\item Let $C_{\geq \max(A)}$ be the set of all elements of $C$ that are bigger than or equal to $\max(A)$.
%\end{enumerate} 
Consider the set $\{C\subseteq [n-1]: A\subseteq C\}$. Let the sign of $C$ to be $(-1)^{|C|}$. 

We first show that  if $n-\max(A)$ is even, then
$$\sum_{A\subseteq C\subseteq[n-1]}(-1)^{|C|}\mathbf{K}_{(C\setminus((C+1)\cup\{1\}),\sigma)}=0$$ using a sign reversing involution $\imath$ with no fixed point defined as follows. 

Let $C_{\geq \max(A)}$ be the set of the elements of $C$ that are bigger than or equal to $\max(A)$. Let 
$$(C_{\geq \max(A)}\setminus \{1\})\cup \{\max(A),n+1\}=\{\max(A)=d_1<d_2<\cdots<d_k=n+1\}.$$ 

The set $\{1\leq i\leq k-1: d_{i+1}-d_{i} \text{~is odd} \}$ is non-empty since $d_{k}-d_1=n+1-\max(A)$ is odd. Let $j=\min(\{1\leq i\leq k-1: d_{i+1}-d_{i} \text{~is odd}\})$; that is $j$  is the smallest integer such that $d_{j+1}-d_j$ is odd. If $ C\cap \{d_j+1, d_j+2,\dots, d_{j+1}\}=\emptyset$, set $m_C=0$, otherwise, let $m_C$ be the largest integer such that 
$$d_{j}+1,d_j+2,\ldots,d_j+m_C \in C\cap \{d_j+1, d_j+2,\dots, d_{j+1}\}.$$
Now define 
$$
\begin{array}{cccc}
\imath: & \{ C\subseteq [n-1]: A\subseteq C\}&\rightarrow &\{ C\subseteq [n-1]: A\subseteq C\}\\
& C& \mapsto & \begin{cases} 
C\setminus \{d_j+m_C\} & \text{if ~} m_C \text{~is odd,}\\
C\cup \{d_j+m_C+1\} & \text{if ~} m_C \text{~is even.}
\end{cases} 
\end{array}
$$
Then $\imath$ is a sign reversing involution with no fixed point. Moreover, for any $C$, we have
$$\imath(C)\setminus((\imath(C)+1)\cup\{1\})=C\setminus((C+1)\cup\{1\}),$$ and so 
$$\mathbf{K}_{(C\setminus((C+1)\cup\{1\}),\sigma)}=\mathbf{K}_{\imath(C)\setminus((\imath(C)+1)\cup\{1\}),\sigma)}.$$
Consequently, 
$$\sum_{A\subseteq C\subseteq[n-1]}(-1)^{|C|}\mathbf{K}_{(C\setminus((C+1)\cup\{1\}),\sigma)}=0.$$

Let $n-\max(A)$ be odd and let $B$ be the unique peak set such that $\odd(B)=\oddd(A)$. We want to show that 
$$\sum_{A\subseteq C\subseteq[n-1]}(-1)^{|C|}\mathbf{K}_{(C\setminus((C+1)\cup\{1\}),\sigma)}=(-1)^{n-1-|B|}\boldeta_{(B,\sigma)}$$ using a sing reversing involution $\jmath$ defined as follows. \\

Define 
$$
\begin{array}{cccc}
\jmath:& \{C\subseteq [n-1]: A\subseteq C\} & \rightarrow& \{C\subseteq [n-1]: A\subseteq C\}\\
& C & \mapsto & 
\begin{cases} 
C & \text{if~}B\subseteq C,\\
 C\setminus\{\ (\min(B)\setminus C)-1 \} & \text{if~}  B\not\subseteq C, (\min(B)\setminus C)-1\in C,\\
 C\cup\{\ (\min(B)\setminus C)-1 \}  & \text{if~}  B\not\subseteq C, (\min(B)\setminus C)-1\not\in C.\\
\end{cases} 
\end{array} 
$$
Note that the set of fixed points of $\jmath$ is 
$$\{C\subseteq [n-1]: A,B\subseteq C\}.$$ 
We have that
$$\{C\setminus ( (C+1)\cup \{1\} ): A,B\subseteq C\subseteq [n-1]\}=\{ D: D\subseteq B \}.$$
The reason for this equality is that if $A,B\subseteq C\subseteq [n-1]$, then 
$$C=[n-1]\setminus \{d-1:d\in D\}$$ for some $D\subseteq B$.
Note that  $|D|=n-1-|C|$. 
Therefore, 
\begin{align*} 
\sum_{A\subseteq C\subseteq[n-1]}(-1)^{|C|}\mathbf{K}_{(C\setminus((C+1)\cup\{1\}),\sigma)}&=\sum_{A,B\subseteq C\subseteq[n-1]}(-1)^{|C|}\mathbf{K}_{(C\setminus((C+1)\cup\{1\}),\sigma)}\\
&=\sum_{D\subseteq B}(-1)^{n-1-|D|}\mathbf{K}_{(D,\sigma)}\\
&=(-1)^{n-1-|B|}\boldeta_{(B,\sigma)}.
\end{align*}

\end{proof} 

\end{document}
